\section{Discussion}
\label{sec:discussion}

\subsection{Implications for benchmarking}

A back-to-back engine benchmark measures the GPU at a clock that a
sequential application never reaches on this device, so engine-level
speedups answer a different question from the one a deployment asks.
The gap is not a constant factor; it depends on the schedule (\cref{tab:ladder}), the arrival rate (\cref{tab:stream}) and
the model's own duty cycle (YOLO26n versus YOLO26s in
\cref{tab:dvfs}). In these experiments it is not even monotone in the
arrival rate. The uint8 and float sequential pipelines have higher
median latency and higher CPU time per frame at 15\,fps than at
30\,fps, which points to a lower CPU clock during longer idle gaps
(\cref{sec:sched}).
This paper therefore suggests that efficiency claims for edge detectors
report at least: (i)~an end-to-end decode-to-detections measurement,
(ii)~the GPU and CPU governors and the power mode, (iii)~the GPU and
CPU clocks logged during the measurement, (iv)~how the harness stores
results, and (v)~the workload scenario (backlog, fixed-rate stream or
single frame), in the spirit of the MLPerf
scenarios~\cite{reddi2020mlperf}. Locking clocks is a useful
diagnostic, but not a neutral benchmarking choice: it removes exactly
the interactions that dominate behavior under the stock governors a
device ships with. Under locked clocks the synchronization call looks
irrelevant on the sequential and prefetch schedules, while under the
default governors it costs a fifth of the throughput. On the two-worker
schedule only the synchronization boundary matters (24\% for stream
synchronization), identically in both regimes (\cref{sec:sched}).
Locking also raises idle power by 1.9\,W, which at camera rates costs
more energy than any pipeline choice saves (\cref{sec:stream}).

\subsection{Guidelines for practitioners}

The ordering that worked on this platform was:
\begin{enumerate}
\item Take host stages off the GPU's critical path first (decode
  prefetch, overlapped post-processing). Choose a wait that does not
  let the CPU governor lower the host's clock: a spinning wait or, where
  a sleeping wait induces frequency collapse, a per-task utilization
  clamp, which needed no root privileges on the test system and costs
  less energy (\cref{sec:sched}). Once the pipeline
  keeps the host busy, the wait policy no longer matters, but the
  synchronization boundary does. In an enqueue-ahead loop, synchronize
  on per-frame events, not on the whole stream (\cref{tab:uclamp2}).
\item Move preprocessing into the engine once the CPU becomes the
  bottleneck. Its value grows from 18\% to $2.4\times$ as the schedule
  improves.
\item Match the schedule to the workload. Use overlap with several
  decode workers for backlogs and high rates. Use the arrival-aware
  variant when frames arrive at camera rates and latency matters. Do
  not lock clocks for a camera application that runs below capacity.
\item Only then evaluate model-level savings (resolution, pruning,
  quantization), in the optimized pipeline and with application-level
  metrics.
\end{enumerate}
Step 4 matters because the pipeline effect is larger than the model
effect. In these measurements, the full YOLO26s pipeline under default
governors is faster and more accurate than the Ultralytics YOLO26n
predictor with locked clocks (\cref{tab:generality}). Model sizes
should be compared within the same, optimized pipeline.

\subsection{Why compression is not useless}

The negative compression results are specific to a 2.6\,M-parameter
detector whose compiled engine takes about 4\,ms, on a device whose
host stages cost more than the engine at full clock. For larger
models whose engine is the bottleneck even in an optimized pipeline
(YOLO26s reaches 95\% of its back-to-back engine rate), reducing engine
work translates directly into throughput, and compression regains its
usual role. The results argue for evaluating compression in the regime
where it will be deployed, not for abandoning it.

\subsection{Relation to frequency-control research}

DVFS research on edge inference usually controls the frequency
directly~\cite{kim2021ztt,han2025dvfs,zhang2026sparsedvfs}. The
present results are complementary: with the stock governor left in
place, the application's own schedule determines the operating point,
and a well-scheduled pipeline lets the governor choose the top clock on
its own. A governor-aware design could go further, for example by
deliberately running at the lowest clock that meets a latency target at
a known frame rate. The 60\,fps streaming results show the potential: the governor ran the optimized pipelines at
405--430\,MHz, and their marginal energy per frame fell below the
backlog value. The CPU results add a second dimension: the CPU and GPU
governors act independently, but through the host pipeline the CPU
clock sets the GPU's duty cycle and therefore its clock
(\cref{sec:sched}). A coordinated policy would have to treat the two as
one control problem. Utilization clamping is a first, coarse step in
this direction: it lets the application tell the CPU governor that its
sleeping thread is on the GPU's critical path.
